\chapter{Breast Biopsy}

\section*{What Is This Test?}
A breast biopsy removes cells or tissue from an area of concern so a pathologist can determine whether it is benign, atypical, in situ, or invasive disease. It is the definitive tissue test for many breast abnormalities.

\section*{Alternative Names}
\begin{itemize}
  \item Breast Tissue Biopsy
  \item Core Needle Biopsy
  \item Stereotactic Breast Biopsy
  \item Ultrasound-Guided Breast Biopsy
\end{itemize}

\section*{Principle}
Imaging localizes the target and guides needle placement. Core biopsy removes tissue cylinders for histology and, when needed, immunohistochemistry or molecular testing. Fine-needle aspiration samples cells but may provide less architectural information.

\section*{Why This Test Is Ordered}
It is performed for a suspicious mammogram, ultrasound, or MRI finding; a new breast or axillary lump; nipple or skin changes; or an abnormality that cannot be safely classified by imaging alone.

\section*{Primary vs Secondary Test}
Biopsy is a secondary diagnostic test after clinical and imaging assessment. Most suspicious findings require image-guided tissue sampling rather than surgical removal as the first step.

\section*{How This Test Is Performed}
After local anesthetic, a needle is guided by ultrasound, mammography, or MRI to obtain several samples. A marker may be placed at the biopsy site, and a mammogram may confirm its position.

\section*{What Preparation Is Needed}
Tell the clinician about anticoagulants, bleeding disorders, allergies, implants, and pregnancy possibility. Wear a two-piece outfit and follow instructions about medicines; do not stop anticoagulants without medical advice.

\section*{Contraindications and Precautions}
Bruising, bleeding, pain, infection, and temporary swelling can occur. A benign result must be concordant with imaging; discordance may require repeat sampling or surgical biopsy.

\section*{Understanding Results}
The pathology report may identify benign changes, atypia, ductal carcinoma in situ, or invasive carcinoma. If cancer is found, grade, hormone receptors, HER2, and other features guide staging and treatment. Results should be reviewed with the breast-care team.

\section*{Follow-up Tests}
Follow-up may include breast imaging, pathology review, receptor testing, breast MRI, lymph-node evaluation, genetic testing, staging studies, or surgical and oncology consultation.

\section*{References}
\begin{enumerate}
  \item MedlinePlus. Breast Biopsy.
  \item American College of Radiology. Practice parameter for the performance of breast biopsy.
  \item National Comprehensive Cancer Network. Breast Cancer Screening and Diagnosis.
\end{enumerate}

\clearpage
\section*{Understanding Results and Follow-up}
The laboratory report should identify the specimen, method, units, reference interval or decision limit, and whether the result is positive, negative, abnormal, or inconclusive. Reference intervals vary with age, sex, pregnancy, specimen type, assay, and laboratory; a result outside a range is not automatically a diagnosis, and a result within a range does not always exclude disease. Discuss unexpected results with the ordering clinician, who can decide whether repeat or confirmatory testing, treatment, or referral is appropriate. Do not change medicines or delay urgent care based on an isolated result.


\section*{Regulatory Status and Authorization Date}
This chapter describes a test category, not a specific commercial product. FDA, CDSCO, EU IVDR, and MHRA approval or registration dates therefore cannot be assigned without the assay manufacturer, product name, instrument, and intended use. For laboratory-developed tests, an individual agency approval date may not exist. See the Preface for the interpretation of regulatory status.

\clearpage
